% Round 16: reconstruction of the historical R6 candidate from the repaired R11 interfaces.
% Status: ACTIVATED after independent root and global-height audits. Historical files unchanged.
\section{Reconstructing the full finite actual trace from independent parameter tori}
\label{r16ft:section}

\paragraph{Provenance and status.}
The multi-parameter diagonal construction was proposed in
\texttt{history\_01\_07/research/round6\_general\_g/envelope\_full\_finite\_actual\_trace\_section.tex}.
It is not a new construction of Round 16. The frozen final register of
Round 14 retains the full trace as a historical candidate and activates
only all-degree PWS trace and the first actual strip. The concrete
boundary-localization gap identified in the initial Round 11 minimal-trace
audit concerned product-moment integration; its final version was repaired
by the refined-skeleton identities PS1--PS6. The proof below restarts from
those repaired interfaces, supplies the multi-parameter steps explicitly,
and does not cite the old full trace as a theorem. Two independent complete
readings, by the root and global-height branches, passed the mathematical
source at SHA256
\texttt{7f1348a590742b822880cbd46cdb71ead9d0a9b4e192913cece1691eb661e694}.
The theorem is activated relative to the explicitly frozen final Round 11
input; the change from that audited source is confined to this status
notice and the matching scope notice below.

Let $K/\mathbf Q_p$ be finite with its fixed absolute value. All varieties
in this section are normal integral separated $K$-varieties, and ordinary
opens in their analytifications carry Mihatsch's actual forms. Every
auxiliary torus is over $K$ and is not a constant-field extension. Write
$D=d'd''$ and $T_s=(\mathbf G_m^s)^{\mathrm{an}}$.

\subsection{The exact frozen input}

We use the final source
\texttt{round11/perfect\_complex/one\_parameter\_trace.tex}, read together
with \texttt{round11/audits/minimal\_trace\_final\_audit.md} and its
Round 12 inherited final audit. Its input assertions are:
\begin{enumerate}
\item The native degree formula for compact actual top forms under a
finite dominant normal map, formula MT2. In particular finite actual
pullback is faithful. This formula precedes any full source trace.
\item Ordinary scalar PL Galois descent, including the explicit orbit
separator and Newton-coefficient construction; the max/min sorting
identities MT3--MT4; and pure PL-curvature trace MT5 with all compact
actual tests and its actual orbit identity.
\item Product-moment integration MT6--MT7, proved through the weighted
product-skeleton identification PS1 and the refined skeletal identities
PS2--PS6. Its input is an arbitrary Euclidean delta coefficient in
$(c,z)$, with a fixed compact set of $v(z)$ values. Its target test is
an arbitrary compact actual form, including additional mixed unit
moments. The map commutes with $d',d''$.
\end{enumerate}
The third item is much more specific than a general proper actual
pushforward. In particular an arbitrary actual form on a product is not
presumed to have such a product-moment presentation.

\subsection{Finitely many independent parameters}

\begin{lemma}[Product-moment integration with $s$ parameters]
\label{r16ft:product}
Let $U$ be an ordinary open in the analytification of a normal integral
$K$-variety, let $c$ be a finite tuple of analytic units on $U$, and let
\[
 \beta=t_{(c,z_1,\ldots,z_s)}^\star\Gamma,
 \qquad \Gamma\in B^{p,q}(\mathbf R^{N+s}),
\]
where the support of $\Gamma$ lies in
$\mathbf R^N\times C$ for a fixed compact $C\subset\mathbf R^s$.
For $p_U:U\times T_s\to U$ and the coordinate projection
$\pi_s:\mathbf R^{N+s}\to\mathbf R^N$, the form
\[
 I^{(s)}_U\beta:=t_c^\star(\pi_{s,*}\Gamma)
 \tag{FT16.1}\label{r16ft:product-output}
\]
is ordinary actual of degree $(p-s,q-s)$ and satisfies
\[
 \int_U I^{(s)}_U\beta\wedge\eta
 =\int_{U\times T_s}\beta\wedge p_U^*\eta
 \tag{FT16.2}\label{r16ft:product-pairing}
\]
for every compact complementary actual test $\eta$. It is independent
of product-moment presentations, commutes with $d',d''$, and is
compatible with ordinary restriction. If $a:V\to U$ is an analytic
map between such ordinary opens, then on these supplied presentations
\[
 a^* I^{(s)}_U\beta
 =I^{(s)}_V (a\times\mathrm{id})^*\beta.
 \tag{FT16.3}\label{r16ft:product-base-change}
\]
\end{lemma}

\begin{proof}
Enclose $C$ in a compact coordinate box $J_1\times\cdots\times J_s$.
Integrate first in $z_s$. With $U\times T_{s-1}$ as base, the input has
one-parameter product moments $(c,z_1,\ldots,z_{s-1};z_s)$.
The frozen one-parameter theorem applies to an arbitrary actual
Euclidean coefficient $\Gamma$ and produces the presentation
\[
 t_{(c,z_1,\ldots,z_{s-1})}^\star(\pi_{z_s,*}\Gamma).
\]
Its remaining parameter support lies in the projection of $C$, hence
in $J_1\times\cdots\times J_{s-1}$. Repeat. Euclidean affine
pushforward composes, so the final expression is
\eqref{r16ft:product-output}.

Here is the compact-test point in that iteration. Choose smooth compact
functions $\theta_j$ on $\mathbf R$ equal to one on neighborhoods of
$J_j$. Before applying the last-parameter pairing, insert
\[
 \left(\prod_{j<s}\theta_j(v(z_j))\right)p_U^*\eta
\]
as test on $U\times T_{s-1}$. This test is actual and compactly
supported: the base support is compact and bounded parameter valuations
lie in compact closed polyannuli. The inserted factors equal one on
the input and output supports. At the next step omit the already
integrated factor and repeat. Each one-parameter equality therefore
uses an allowed compact actual test, and the cutoffs disappear from
the final equality. No unbounded pullback test is used.

The one-parameter proof identifies every refined base skeleton with
its Gauss product, with base weight times parameter weight one. At
each induction step its base includes all the preceding parameter
coordinates and every additional test moment. Thus the native
calibration and all-test scope persist, rather than merely a comparison
on product test tensors. The argument never replaces an ordinary
compact support by an assertion that its tropical image avoids every
chart-boundary image; that issue is handled by the frozen PS2--PS6
skeletal identities at each step.

If two product presentations give the same actual input, the right
side of \eqref{r16ft:product-pairing} is the same, so current
injectivity makes the outputs identical. The same argument proves
ordinary compatibility. Derivatives preserve the parameter support;
Euclidean affine differentiation and the actual moment pullback give
the asserted commutation.

For \eqref{r16ft:product-base-change}, the pulled-back input is the
supplied presentation $t_{(c\circ a,z)}^\star\Gamma$. Both displayed
outputs are exactly $t_{c\circ a}^\star(\pi_{s,*}\Gamma)$. This is
functoriality of actual pullback for a fixed presentation, not a
base-change theorem for arbitrary currents or proper actual outputs.
\end{proof}

\begin{lemma}[Several parameters in pure-curvature trace]
\label{r16ft:uniform}
Let $q:V\to U$ be the restriction of a finite dominant normal algebraic
map. Suppose finitely many PL scalar functions $u_i$ on $V\times T_s$
are jointly finite PL in a fixed source-unit tuple $f$ and the
independent parameter valuations $t=(v(z_1),\ldots,v(z_s))$.
Locally around every whole finite base fiber, the pure-curvature trace
of each $\prod_iDu_i$ has a presentation
\[
 t_{(c,z)}^\star K_0
 \tag{FT16.4}\label{r16ft:uniform-output}
\]
for a single finite tuple of downstairs analytic units $c$, independent
of $t$, and a Euclidean delta-form $K_0$. The base shrinking can be
chosen independently of $t$. The trace retains its native degree and
its all-actual-test pairing on $U\times T_s$.
\end{lemma}

\begin{proof}
It suffices first to consider a Galois map with group $G$. Enlarge $f$
by its $G$-orbit, so $G$ acts by coordinate permutation on its logarithm
space $\mathbf R^a$. The frozen scalar descent uses the proper finite
PL orbit separator
\[
 J:\mathbf R^a\longrightarrow\mathbf R^b
\]
formed from ordered orbit multisets in finitely many integral
Vandermonde directions. For parameter-dependent invariant PL scalars
one uses $J\times\mathrm{id}_{\mathbf R^s}$. This map still has
precisely the $G$-orbits as fibers and is proper. On a finite common
polyhedral subdivision it has affine inverses onto its image cells.
Invariance identifies the expressions on overlapping image cells.
A PL ambient extension gives a finite PL expression in $(J\ell,t)$.
When the initial expression is invariant only on the realized source
image, average its ambient expression over $G$ first; this changes
neither its actual restriction nor its finite PL character.

The tuple expressing $J\ell$ downstairs is constructed exclusively
from orbit polynomials of source monomials in $f$. None of those
monomials involves $z_j$. Their invariant coefficients are analytic
functions downstairs. At the fixed base point keep the nonzero
coefficients as units; shrink so coefficients vanishing there stay
strictly above the endpoint chord of the corresponding Newton polygon.
The constant coefficient is a unit and the leading coefficient is one,
so that exclusion is uniform. Exactly as in the frozen scalar proof,
ordered root valuations are therefore finite PL functions of a fixed
finite downstairs unit tuple $c$. The choice of coefficients and the
shrinking of $U$ do not depend on any $t_j$.

Sort the $G$-orbit of any one potential, now viewed as a PL function
of $(\ell,t)$. Max/min sorting preserves the sums of every power of
its curvature by MT3. Each ordered potential is invariant, hence is a
PL scalar in $(v(c),t)$ by the preceding paragraph. Apply polarization
to treat mixed products $Du_1\cdots Du_k$. Differentiation is performed
only after scalar descent and uses genuine actual pullback in the
unit tuple $(c,z)$. This constructs \eqref{r16ft:uniform-output}; it
never pulls a delta-form through the nonlinear Newton map.

The same sorting proof gives the actual Galois orbit identity on the
whole product, including all walls of the descended PL potentials.
Its pairing with an arbitrary compact actual product-space test is
proved by the native top-degree formula for $q\times\mathrm{id}$,
whose degree is $|G|$, and invariance under the source automorphisms.
The test may have arbitrary mixed analytic unit moments. Thus no
product-test restriction has entered this step.

For a non-Galois map take a normal Galois closure, pull the finite
source-unit tuple to it, and divide the Galois trace by the full degree
of that closure over $V$. The native degree formula proves the same
pairing. All orbit coefficients still involve source units alone,
so the parameter-uniform presentation remains valid.
\end{proof}

\subsection{The exact diagonal evaluates arbitrary delta coefficients}

\begin{lemma}[Actual graph evaluation]
\label{r16ft:graph}
Let $f=(f_1,\ldots,f_s)$ be analytic units on an ordinary open $V$,
with $v(f)(V)$ contained in a compact subset of $\mathbf R^s$, and let
$\gamma\in B^{p,q}(\mathbf R^s)$. Choose
$\chi\in C_c^\infty(\mathbf R^s)$ equal to one on a neighborhood of
that subset. On $V\times T_s$ put
\[
 \Gamma_f=\prod_{i=1}^sD\max\{v(f_i)-v(z_i),0\}.
\]
Then the input below is a supplied product-moment input and
\[
 I^{(s)}_V\bigl(\chi(t)\gamma(t)\wedge\Gamma_f\bigr)
                      =t_f^\star\gamma
 \tag{FT16.5}\label{r16ft:graph-evaluation}
\]
as ordinary actual forms, hence with the full pairing of
Lemma~\ref{r16ft:product}.
\end{lemma}

\begin{proof}
In $\mathbf R^s_u\times\mathbf R^s_t$ the product
\[
 \Delta=\prod_{i=1}^sD\max\{u_i-t_i,0\}
\]
is the diagonal cycle with its standard weight. Its defining affine
normal coordinates are $u_i-t_i$, so each normal determinant is one.
The Euclidean delta intersection identity is
\[
 \pi_{u,*}\bigl(\pi_t^*(\chi\gamma)\wedge\Delta\bigr)
                         =\chi\gamma.
 \tag{FT16.6}\label{r16ft:euclidean-diagonal}
\]
This is the identity-graph pullback formula in Mihatsch's Euclidean
intersection theory and applies to every Euclidean delta-form
$\gamma$, not only smooth coefficients. In particular it is an
intersection product of delta-forms, not multiplication or restriction
of arbitrary distributions.

Pull \eqref{r16ft:euclidean-diagonal} through the genuine unit tuple
$f$ and use the definition \eqref{r16ft:product-output}. On the source,
$\chi(v(f))=1$ on an ordinary neighborhood of its entire realized
moment image, so $t_f^\star(\chi\gamma)=t_f^\star\gamma$.
The product-moment theorem proves the corresponding identity on all
refined skeletons and against every actual compact target test. The
graph kernel has even total degree, so no sign is introduced by the
written order of its factors.
\end{proof}

\subsection{The trace construction}

\begin{theorem}[Full finite actual trace over the original field]
\label{r16ft:full-trace}
Let $q:Y\to X$ be finite dominant between normal integral separated
$K$-varieties, of full generic degree $d$. There is a canonical map of
ordinary sheaves in every bidegree
\[
 \operatorname{Tr}_q:q_*B_Y^{p,q}\longrightarrow B_X^{p,q}
 \tag{FT16.7}\label{r16ft:trace-map}
\]
characterized by
\[
 \int_X\operatorname{Tr}_q\omega\wedge\eta
       =\int_Y\omega\wedge q^*\eta
 \tag{FT16.8}\label{r16ft:trace-pairing}
\]
for every compact complementary actual test on each ordinary open.
It commutes with $d',d''$, is transitive, satisfies the projection
formula for every actual downstairs factor, and
\[
             \operatorname{Tr}_q(q^*\alpha)=d\alpha.
 \tag{FT16.9}\label{r16ft:degree}
\]
The original field and native residue and constant-field degrees are
retained.
\end{theorem}

\begin{proof}
Fix $x\in X^{\mathrm{an}}$. Its inverse image is a finite set. Around
each of its points choose disjoint ordinary presentation neighborhoods
for $\omega$, and then smaller relatively compact neighborhoods with
unit logarithms in compact subsets of the respective Euclidean
presentation domains. Properness of $q$ shrinks an ordinary neighborhood
$U$ of $x$ so that $V=q^{-1}(U)$ lies in their union. The resulting
source pieces are open and closed in $V$.

Combine the finitely many source presentations as follows. Extend each
unit from its own clopen piece by the constant one on the other pieces.
Add a marker unit taking distinct constants $\varpi^{a_i}\in K^\times$
on the pieces, for a fixed $\varpi$ with nonzero valuation. Each local
Euclidean presenter can be multiplied by a smooth cutoff compactly
supported in its original Euclidean open, equal to one near its compact
realized image, and then extended by zero. A smooth function of the
marker coordinate selects each piece. Summing the pullbacks of those
extended presenters under coordinate projections therefore gives one
finite tuple $f$ and one genuine Euclidean delta-form $\gamma$ such that
\[
                     \omega=t_f^\star\gamma\quad\hbox{on }V.
 \tag{FT16.10}\label{r16ft:whole-fiber}
\]
The source log image remains bounded after this shrinking. No extension
from an intrinsic tropical image or closed ambient lift is used.

Let $s$ be the length of $f$. Form $\Gamma_f$ as in
Lemma~\ref{r16ft:graph} and put $q_T=q\times\mathrm{id}_{T_s}$.
It is a product of PL curvatures, so the independently constructed
pure-curvature trace gives
\[
                    K_q=\operatorname{Tr}^{\mathrm{curv}}_{q_T}\Gamma_f.
\]
Lemma~\ref{r16ft:uniform} supplies a single product-moment expression
$K_q=t_{(c,z)}^\star K_0$, with base units and shrinking independent of
the parameter. Choose $\chi$ equal to one near $v(f)(V)$. Then
\[
 \Theta=\chi(t)\gamma(t)\wedge K_q
       =t_{(c,z)}^\star\bigl(\pi_t^*(\chi\gamma)\wedge K_0\bigr).
 \tag{FT16.11}\label{r16ft:integrand}
\]
has compact parameter support and hence lies in the proved
product-moment integration domain. Define
\[
           \operatorname{Tr}_q\omega:=I^{(s)}_U\Theta.
 \tag{FT16.12}\label{r16ft:trace-construction}
\]
Its degree is $(p+s-s,q+s-s)=(p,q)$.

For a compact complementary actual test $\eta$ on $U$, the three
independently established pairings give
\[
\begin{split}
 \int_U I^{(s)}_U\Theta\wedge\eta
 &=\int_{U\times T_s}\chi(t)\gamma(t)\wedge K_q\wedge p_U^*\eta\\
 &=\int_{V\times T_s}\chi(t)\gamma(t)\wedge\Gamma_f
                                      \wedge p_V^*q^*\eta\\
 &=\int_V\omega\wedge q^*\eta.
\end{split}
 \tag{FT16.13}\label{r16ft:three-pairings}
\]
The first is Lemma~\ref{r16ft:product}; the second is the
pure-curvature trace with the actual downstairs multiplier
$\chi(t)\gamma(t)\wedge p_U^*\eta$; the third is
Lemma~\ref{r16ft:graph} and its all-test product pairing.
Every test in the middle step is compact: it has compact base support
and compact logarithmic parameter support. Its degree is complementary
to the graph kernel. The last target test $q^*\eta$ is compact because
$q$ is finite. Thus all three equalities apply to arbitrary actual
$\eta$, not merely smooth or product tests.

If a different presenter, cutoff, marker tuple, Galois closure or
ordinary neighborhood is chosen, the right side of
\eqref{r16ft:three-pairings} is unchanged. The outputs are actual,
so their equality follows from Mihatsch's current injectivity. In
particular a source presentation of the zero actual form produces
zero. These equalities on ordinary overlaps glue the sheaf map.
Actual Stokes gives commutation with both differentials; since the
construction exists in every degree, there is no degree-range proviso.
Apply the pairing twice to obtain transitivity. Apply it with a
product of any actual downstairs factor and a complementary test to
obtain the projection formula. Finally the native degree formula on
$q^*(\alpha\wedge\eta)$ gives \eqref{r16ft:degree}. No auxiliary field
has been substituted for $K$ at any stage.
\end{proof}

\begin{corollary}[Orbit identity and all six differential operators]
\label{r16ft:orbit}
For finite Galois $q$ with group $G$,
\[
 q^*\operatorname{Tr}_q\omega=\sum_{g\in G}g^*\omega.
 \tag{FT16.14}\label{r16ft:orbit-formula}
\]
For all finite dominant normal maps the trace preserves the intrinsic
actual trigrading and commutes with
$d',d'',d'_P,d''_P,\partial',\partial''$.
In particular every invariant actual form descends through a finite
Galois quotient as an ordinary actual form, in every bidegree.
\end{corollary}

\begin{proof}
In the local construction, pure-curvature trace has the orbit identity
$q_T^*K_q=\sum_g g^*\Gamma_f$. Pull back
\eqref{r16ft:trace-construction} using the supplied-presentation
base-change identity \eqref{r16ft:product-base-change}. The multiplier
$\chi(t)\gamma(t)$ is fixed under $G$, and graph evaluation for the
source tuple $g^*f$ gives \eqref{r16ft:orbit-formula}. The cutoff is
one near all source values since $g$ permutes the whole source domain.

Analytic pullback preserves the intrinsic trigrading and all six
operators. The orbit identity and faithful finite pullback imply
trigrading preservation and differential commutation for a Galois map.
For a general finite map choose a normal Galois closure $p:W\to Y$,
$q'=qp$. The defining all-test pairings and the native top-degree
formula imply
\[
 \operatorname{Tr}_q\omega
      =[K(W):K(Y)]^{-1}\operatorname{Tr}_{q'}(p^*\omega).
\]
This gives the assertions in general. For an invariant form, divide its
trace by $|G|$ in \eqref{r16ft:orbit-formula}; uniqueness again follows
from faithful pullback. The degree contains every constant-field and
residual degree, not merely the number of topological fiber points.
\end{proof}

\begin{corollary}[Finite cohomological detection in all degrees]
\label{r16ft:cohomology}
Finite dominant normal pullback is injective on actual Bott--Chern,
Aeppli, $d'$ and $d''$ cohomology in every bidegree. If
$q^*\Omega=D\beta$ for an actual primitive of any degree, then
\[
          \Psi=d^{-1}\operatorname{Tr}_q\beta,
                     \qquad D\Psi=\Omega.
 \tag{FT16.15}\label{r16ft:primitive}
\]
\end{corollary}

\begin{proof}
The normalized trace is a left inverse to pullback and commutes with
both differentials. It therefore preserves each of the stated
numerators and boundary spaces. The displayed primitive formula is
an equality of actual forms and retains all compact actual tests.
\end{proof}

\paragraph{Exact scope.}
Normality of the target is retained; the normalization of a nodal curve
already fails to admit trace for all ordinary scalar forms. The theorem
does not define a proper actual pushforward for arbitrary analytic maps.
It gives no new PWS replacement of higher-degree primitives, and the
Round 15 trigrading counterexample remains valid. The activated trace does
remove the finite-descent boundary in the Round 15 actual projective-line
projection: UG and AD already construct the upstairs output in every
degree, and the present full trace can descend that output. The latter
projection retains the smooth-mixed-test scope of UG/AD unless a separate
all-actual-test proof is supplied. Neither finite trace nor that relative
projection alone proves the homological vanishing needed for general
strong G or the global Beilinson--Bloch height.
